% Options for packages loaded elsewhere
\PassOptionsToPackage{unicode}{hyperref}
\PassOptionsToPackage{hyphens}{url}
\PassOptionsToPackage{dvipsnames,svgnames,x11names}{xcolor}
\documentclass[
  11pt,
  a4paper,
]{article}
\usepackage{xcolor}
\usepackage[margin=2.6cm]{geometry}
\usepackage{amsmath,amssymb}
\setcounter{secnumdepth}{5}
\usepackage{iftex}
\ifPDFTeX
  \usepackage[T1]{fontenc}
  \usepackage[utf8]{inputenc}
  \usepackage{textcomp} % provide euro and other symbols
\else % if luatex or xetex
  \usepackage{unicode-math} % this also loads fontspec
  \defaultfontfeatures{Scale=MatchLowercase}
  \defaultfontfeatures[\rmfamily]{Ligatures=TeX,Scale=1}
\fi
\usepackage{lmodern}
\ifPDFTeX\else
  % xetex/luatex font selection
  \setmainfont[]{Libertinus Serif}
  \setmonofont[Scale=0.85]{Latin Modern Mono}
  \setmathfont[]{Latin Modern Math}
\fi
% Use upquote if available, for straight quotes in verbatim environments
\IfFileExists{upquote.sty}{\usepackage{upquote}}{}
\IfFileExists{microtype.sty}{% use microtype if available
  \usepackage[]{microtype}
  \UseMicrotypeSet[protrusion]{basicmath} % disable protrusion for tt fonts
}{}
\makeatletter
\@ifundefined{KOMAClassName}{% if non-KOMA class
  \IfFileExists{parskip.sty}{%
    \usepackage{parskip}
  }{% else
    \setlength{\parindent}{0pt}
    \setlength{\parskip}{6pt plus 2pt minus 1pt}}
}{% if KOMA class
  \KOMAoptions{parskip=half}}
\makeatother
\usepackage{longtable,booktabs,array}
\newcounter{none} % for unnumbered tables
\usepackage{calc} % for calculating minipage widths
% Correct order of tables after \paragraph or \subparagraph
\usepackage{etoolbox}
\makeatletter
\patchcmd\longtable{\par}{\if@noskipsec\mbox{}\fi\par}{}{}
\makeatother
% Allow footnotes in longtable head/foot
\IfFileExists{footnotehyper.sty}{\usepackage{footnotehyper}}{\usepackage{footnote}}
\makesavenoteenv{longtable}
\setlength{\emergencystretch}{3em} % prevent overfull lines
\providecommand{\tightlist}{%
  \setlength{\itemsep}{0pt}\setlength{\parskip}{0pt}}

\providecommand{\E}{\mathbb{E}}
\providecommand{\N}{\mathbb{N}}
\providecommand{\Z}{\mathbb{Z}}
\providecommand{\R}{\mathbb{R}}
\providecommand{\eps}{\varepsilon}
\providecommand{\ceil}[1]{\left\lceil #1\right\rceil}
\providecommand{\floor}[1]{\left\lfloor #1\right\rfloor}
\AtBeginDocument{\let\setminus\smallsetminus}
\usepackage[most]{tcolorbox}
\newtcolorbox{warning}{colback=red!5,colframe=red!65!black,boxrule=0.8pt,arc=2pt,left=6pt,right=6pt,top=5pt,bottom=5pt}
\usepackage{etoolbox}
\AtBeginEnvironment{longtable}{\small}
\setlength{\emergencystretch}{3em}
\usepackage{fancyhdr}
\pagestyle{fancy}
\fancyhf{}
\fancyhead[L]{\small\itshape Candidate result 2604.09449\_\_03 --- GPT-6 Astra, awaiting review}
\fancyhead[R]{\small\thepage}
\renewcommand{\headrulewidth}{0.2pt}
\usepackage{bookmark}
\IfFileExists{xurl.sty}{\usepackage{xurl}}{} % add URL line breaks if available
\urlstyle{same}
\hypersetup{
  pdftitle={A candidate proof of Problem 6.4 in Colour-balanced subgraphs},
  pdfauthor={github.com/graph-theory-AI --- machine-generated candidate writeup},
  colorlinks=true,
  linkcolor={blue!50!black},
  filecolor={Maroon},
  citecolor={Blue},
  urlcolor={blue!50!black},
  pdfcreator={LaTeX via pandoc}}

\title{A candidate proof of Problem 6.4 in Colour-balanced subgraphs}
\author{\href{https://github.com/graph-theory-AI}{github.com/graph-theory-AI}
--- machine-generated candidate writeup}
\date{17 September 2026}

\begin{document}
\maketitle
\begin{abstract}
The requested O(k\^{}2) bound improves to O(k log(k+1)) for Hamilton
cycles in both complete and complete bipartite graphs, including for
bounded vector labels in an arbitrary norm. This candidate proof was
generated by GPT-6 Astra in a single model pass for the Graph Theory AI
project. It was selected from the campaign because the model marked it
\texttt{would\_publish:\ true}; it has not been checked by the
adversarial referee, a human mathematician, or a novelty search.
\end{abstract}

\section{The problem and the claimed
result}\label{the-problem-and-the-claimed-result}

This document concerns
\href{https://graph-theory-ai.github.io/graph-conjectures/arxiv/2604.09449__03/}{Problem
6.4} from \emph{Colour-balanced subgraphs}.

\subsection{Problem statement}\label{problem-statement}

Problem. Improve the bounds on \(f_{h}(H)\) for a Hamilton cycle \(H\)
in complete or complete bipartite graphs.

\subsection{Catalog context}\label{catalog-context}

The general upper bound from Theorem 1.2 gives \(O(k^2)\) for Hamilton
cycles in complete graphs, and the same bound holds for complete
bipartite graphs by the paper's methods. A necklace-splitting argument
shows any improvement to the bipartite matching bound yields an equal
improvement here, so progress on either problem benefits the other.

\subsection{Claimed result}\label{claimed-result}

\textbf{Model verdict: proved} (confidence: high).

The requested O(k\^{}2) bound improves to O(k log(k+1)) for Hamilton
cycles in both complete and complete bipartite graphs, including for
bounded vector labels in an arbitrary norm.

\emph{The model's own caveat:} Optimality and priority relative to later
revisions of the source paper are not established.

\section{Model writeup}\label{model-writeup}

The text below is the model output verbatim, apart from moving its
machine-readable verdict block into the summary above and shifting
Markdown heading levels for this document. Mechanical TeX defects and
equation tags were normalized where needed for compilation.

\subsection{\texorpdfstring{An \(O(k\log k)\)
bound}{An O(k\textbackslash log k) bound}}\label{an-oklog-k-bound}

The earlier attempt's midpoint-splicing and approximate-convexity
arguments are sound. Below I verify both, sharpen the midpoint estimate,
and extend the complete-graph argument to arbitrary bounded vector
labels. Thus the improvement does not depend on interpreting \(f_h\)
only as a recolouring statistic.

Here ``proved'' refers to the requested improvement in the upper bound,
not to determining its optimal order.

\subsubsection{1. Precise statement}\label{precise-statement}

For \(d\ge 1\), put \[
B_d=\lceil\log_2 d\rceil+2.
\] If \(h\) labels the edges of a graph by vectors, write \[
h(F)=\sum_{e\in E(F)}h(e).
\]

\paragraph{Theorem 1}\label{theorem-1}

Let \(X\) be a real \(d\)-dimensional normed space, and suppose that \[
h:E(G)\longrightarrow X,\qquad \|h(e)\|\le L.
\]

\begin{enumerate}
\def\labelenumi{\arabic{enumi}.}
\item
  If \(G=K_{n,n}\), where \(n\ge2\), then \(G\) has a Hamilton cycle
  \(H\) satisfying \[
  \left\|h(H)-\frac2n h(G)\right\|
  \le 6LdB_d.
  \qquad\text{(1)}
  \]
\item
  If \(G=K_N\), where \(N\ge3\), then \(G\) has a Hamilton cycle \(H\)
  satisfying \[
  \left\|h(H)-\frac{2}{N-1}h(G)\right\|
  \le 6LdB_d+8Ld+2L.
  \qquad\text{(2)}
  \]
\end{enumerate}

The centring vectors in (1) and (2) are the expected label sums of
uniformly random Hamilton cycles. In particular, for unit-bounded,
zero-sum \(d\)-dimensional labels, the least norm of a Hamilton-cycle
label sum is \(O(d\log(d+1))\).

For the recolouring formulation, the theorem gives the following
explicit bounds.

\paragraph{Corollary 2}\label{corollary-2}

Let \(k\ge2\), and suppose that the host graph is \(k\)-edge-coloured
with equally large colour classes.

\begin{itemize}
\item
  If the host is \(K_{n,n}\), \(n\ge2\), and \(k\mid 2n\), some Hamilton
  cycle can be made colour-balanced by recolouring at most \[
  3kB_k
  \qquad\text{(3)}
  \] of its edges.
\item
  If the host is \(K_N\), \(N\ge3\), and \(k\mid N\), some Hamilton
  cycle can be made colour-balanced by recolouring at most \[
  3kB_k+4k+1
  \qquad\text{(4)}
  \] of its edges.
\end{itemize}

Consequently, the quadratic upper bound in the question is replaced by
\[
\boxed{O(k\log(k+1)).}
\]

The substantive ingredient is an improved vector-labelled bipartite
matching bound.

\begin{center}\rule{0.5\linewidth}{0.5pt}\end{center}

\subsubsection{2. Halving a vector-valued
necklace}\label{halving-a-vector-valued-necklace}

We first establish the continuous splitting fact needed in the matching
argument.

\paragraph{Lemma 3}\label{lemma-3}

Let \(f:[0,r]\to\mathbb R^d\) be integrable. There is a set
\(A\subseteq[0,r]\) whose indicator is constant between at most \(d\)
cut points and such that \[
\int_A f(x)\,dx=\frac12\int_0^r f(x)\,dx.
\qquad\text{(5)}
\]

\subparagraph{Proof}\label{proof}

For \(x=(x_0,\ldots,x_d)\in S^d\), partition \([0,r]\) into consecutive
intervals \(J_0,\ldots,J_d\) of lengths \[
r x_0^2,\ldots,r x_d^2.
\] Define \[
F(x)=\sum_{i=0}^d \operatorname{sgn}(x_i)\int_{J_i}f(t)\,dt.
\] The map \(F:S^d\to\mathbb R^d\) is continuous: when a coordinate
tends to zero, its interval has vanishing length, so the corresponding
integral tends to zero. It is also odd.

By the Borsuk--Ulam theorem, \(F\) has a zero. Taking \(A\) to be the
union of the positively signed intervals gives (5). There are at most
\(d\) internal cut points. \(\square\)

No positivity assumption on the coordinates of \(f\) is needed.

\begin{center}\rule{0.5\linewidth}{0.5pt}\end{center}

\subsubsection{3. Approximate midpoint closure for matching
sums}\label{approximate-midpoint-closure-for-matching-sums}

\paragraph{Lemma 4}\label{lemma-4}

Suppose that the edges of \(K_{r,r}\) are labelled by vectors in \(X\),
each of norm at most \(L\). For any two perfect matchings \(P,Q\), there
is a perfect matching \(R\) such that \[
\left\|h(R)-\frac{h(P)+h(Q)}2\right\|\le 3Ld.
\qquad\text{(6)}
\]

\subparagraph{Proof}\label{proof-1}

Let the bipartition be \(U\cup V\). Regard \(P,Q\) as bijections from
\(V\) to \(U\), and define a permutation \(\sigma\) of \(V\) by \[
Q(v)=P(\sigma(v)).
\] Order the vertices of \(V\) by concatenating cyclic orders of the
cycles of \(\sigma\). Write the resulting order as \(v_1,\ldots,v_r\).
Let \(P_i,Q_i\) denote the respective matching edges incident with
\(v_i\).

On \(I_i=[i-1,i]\), put \[
f(x)=h(P_i)-h(Q_i).
\] Apply Lemma 3, after identifying \(X\) linearly with \(\mathbb R^d\),
and set \[
\alpha_i=|A\cap I_i|.
\] Then \[
\sum_i\alpha_i\bigl(h(P_i)-h(Q_i)\bigr)
=\frac12\sum_i\bigl(h(P_i)-h(Q_i)\bigr).
\qquad\text{(7)}
\]

Choose \(s_i\in\{0,1\}\) nearest to \(\alpha_i\), breaking ties
arbitrarily. At most \(d\) of the numbers \(\alpha_i\) lie strictly
between zero and one, since each such unit interval contains an internal
cut point. Therefore \[
\sum_i|s_i-\alpha_i|\le \frac d2.
\qquad\text{(8)}
\]

Select \(P_i\) if \(s_i=1\), and \(Q_i\) if \(s_i=0\). This gives one
selected edge at every vertex of \(V\), though it need not yet be a
matching. Let \(z\) be its label sum. By (7) and (8), \begin{equation*}
\begin{aligned}
\left\|z-\frac{h(P)+h(Q)}2\right\|
&=
\left\|
\sum_i(s_i-\alpha_i)\bigl(h(P_i)-h(Q_i)\bigr)
\right\|\\
&\le 2L\sum_i|s_i-\alpha_i|
\le Ld.
\end{aligned}
\qquad\text{(9)}
\end{equation*}

We now bound the number of matching defects.

For each \(i\), the chosen value \(s_i\) occurs as the indicator of
\(A\) at some non-cut point in the interior of \(I_i\). Choosing such
points in increasing order shows that the linear binary sequence \[
s_1,\ldots,s_r
\] has at most \(d\) changes.

Consider one cycle block of \(\sigma\), and let \(c\) be its number of
internal linear changes. Closing that block cyclically adds at most one
change; if it adds one, then \(c\ge1\). Hence the number of cyclic
changes in that block is at most \(2c\). Summing over all blocks, the
total number \(b\) of cyclic changes satisfies \[
b\le 2d.
\qquad\text{(10)}
\]

For a cycle written so that \(\sigma(v_i)=v_{i+1}\), the degree of
\(P(v_i)\) in the selected edge set is \[
1+s_i-s_{i-1},
\] with cyclic subscripts. Consequently, there are exactly \(b/2\)
vertices of \(U\) of degree two and \(b/2\) of degree zero.

Delete one selected edge at each degree-two vertex. This frees \(b/2\)
vertices of \(V\). Match them arbitrarily to the degree-zero vertices of
\(U\), using completeness of the host. The resulting graph is a perfect
matching \(R\), obtained by replacing \(b/2\le d\) selected edges. Thus
\[
\|h(R)-z\|\le 2L\frac b2\le 2Ld.
\qquad\text{(11)}
\] Combining (9) and (11) proves (6). \(\square\)

The important feature is that the error is independent of \(r\).

\begin{center}\rule{0.5\linewidth}{0.5pt}\end{center}

\subsubsection{4. Approximate convexification costs only a
logarithm}\label{approximate-convexification-costs-only-a-logarithm}

\paragraph{Lemma 5}\label{lemma-5}

Let \(S\) be a finite subset of a \(d\)-dimensional normed space.
Suppose that for every \(u,v\in S\), some \(w\in S\) satisfies \[
\left\|w-\frac{u+v}{2}\right\|\le D.
\qquad\text{(12)}
\] Then every \(x\in\operatorname{conv}(S)\) satisfies \[
\operatorname{dist}(x,S)\le DB_d.
\qquad\text{(13)}
\]

\subparagraph{Proof}\label{proof-2}

By Carathéodory's theorem, write \[
x=\sum_{i=1}^m\beta_i x_i,
\qquad
x_i\in S,\quad m\le d+1.
\] The case \(m=1\) is immediate.

First suppose that all coefficients have a common dyadic denominator
\(2^q\). Form a complete binary tree with \(2^q\) leaves, placing the
required copies of \(x_1,\ldots,x_m\) in consecutive blocks.

Assign representatives in \(S\) to the internal nodes, proceeding
upwards. If all leaves below a node are equal, choose that common point.
Otherwise use (12) on its two child representatives.

At depth \(j\) from the root, the number \(H_j\) of nodes whose leaves
are not all equal is at most \[
H_j\le \min(2^j,m-1).
\qquad\text{(14)}
\] Indeed, their disjoint leaf intervals must each contain a boundary
between consecutive blocks.

An error introduced at depth \(j\) contributes to the root with
coefficient \(2^{-j}\). Therefore the root representative \(z\in S\)
satisfies \begin{equation*}
\begin{aligned}
\|z-x\|
&\le D\sum_{j=0}^{q-1}2^{-j}H_j\\
&\le D\sum_{j\ge0}\min\bigl(1,(m-1)2^{-j}\bigr)\\
&\le D\bigl(\lceil\log_2(m-1)\rceil+2\bigr)\\
&\le DB_d.
\end{aligned}
\qquad\text{(15)}
\end{equation*}

For arbitrary coefficients, approximate them by dyadic coefficients
summing to one. Since \(S\) is finite, a subsequence of the resulting
root representatives is constant. Passing to the limit proves (13).
\(\square\)

This is where the quadratic dependence is avoided: a sequential
convexification can pay an \(O(d)\) error \(O(d)\) times, whereas the
weighted dyadic construction pays only \(O(\log d)\) levels.

\paragraph{Proposition 6: vector matching
bound}\label{proposition-6-vector-matching-bound}

For an edge-labelling of \(K_{r,r}\) in \(X\) with \(\|h(e)\|\le L\),
there is a perfect matching \(M\) such that \[
\left\|h(M)-\frac1r h(K_{r,r})\right\|
\le 3LdB_d.
\qquad\text{(16)}
\]

In fact, every point in the convex hull of the perfect-matching label
sums is within \(3LdB_d\) of a perfect-matching label sum.

\subparagraph{Proof}\label{proof-3}

Apply Lemma 5 to \[
S=\{h(M):M\text{ is a perfect matching}\},
\] using Lemma 4 with \(D=3Ld\).

A uniformly random perfect matching contains each edge with probability
\(1/r\), so \[
\frac1r h(K_{r,r})\in\operatorname{conv}(S).
\] This proves (16). \(\square\)

\begin{center}\rule{0.5\linewidth}{0.5pt}\end{center}

\subsubsection{5. Hamilton cycles in complete bipartite
graphs}\label{hamilton-cycles-in-complete-bipartite-graphs}

Let the parts of \(K_{n,n}\) be \[
A=\{a_1,\ldots,a_n\},\qquad B,
\] with \(A\) cyclically ordered and \(a_{n+1}=a_1\).

Construct an auxiliary \(K_{n,n}\) between the gaps \[
g_i=(a_i,a_{i+1})
\] and the vertices of \(B\). Give its edges the labels \[
\lambda(g_i b)=h(a_i b)+h(a_{i+1}b).
\qquad\text{(17)}
\] These labels have norm at most \(2L\).

A perfect matching in the auxiliary graph assigns distinct vertices
\(b_i\) to the gaps and produces the Hamilton cycle \[
a_1b_1a_2b_2\cdots a_nb_na_1.
\qquad\text{(18)}
\] Its \(h\)-sum is exactly the auxiliary matching's \(\lambda\)-sum.

Each original edge appears in exactly two auxiliary labels. Hence \[
\frac1n\sum_{i,b}\lambda(g_i b)=\frac2n h(K_{n,n}).
\qquad\text{(19)}
\] Applying Proposition 6 with label bound \(2L\) gives (1).

This construction also works for \(n=2\); it produces the unique
four-cycle.

\begin{center}\rule{0.5\linewidth}{0.5pt}\end{center}

\subsubsection{6. Hamilton cycles in complete
graphs}\label{hamilton-cycles-in-complete-graphs}

We need a nearly balanced cut, now for arbitrary normed vector labels.

\paragraph{Lemma 7}\label{lemma-7}

Let \(N\ge4\), and suppose the edges of \(K_N\) have labels in \(X\) of
norm at most \(L\). Put \[
a=\left\lceil\frac N2\right\rceil,\qquad
b=\left\lfloor\frac N2\right\rfloor,\qquad
p=\frac{2ab}{N(N-1)}.
\] There is a partition \(V(K_N)=A\cup B\), with \(|A|=a\), \(|B|=b\),
such that \[
\left\|h(E(A,B))-p\,h(K_N)\right\|\le 2LdN.
\qquad\text{(20)}
\]

\subparagraph{Proof}\label{proof-4}

Choose \(A\) uniformly among the \(a\)-subsets, and let \(I_e\) indicate
that \(e\) crosses the partition. Thus \(\mathbb E I_e=p\).

We first record a scalar variance bound. Let \(|w_e|\le L\), and put \[
Z=\sum_e w_e(I_e-p).
\] For adjacent distinct edges \(e,f\), \[
\mathbb P(I_e=I_f=1)=\frac p2,
\] and therefore \[
\left|\operatorname{Cov}(I_e,I_f)\right|
=\left|p\left(\frac12-p\right)\right|
\le \frac1N.
\qquad\text{(21)}
\] For disjoint edges, direct enumeration gives \[
\operatorname{Cov}(I_e,I_f)=
\begin{cases}
\dfrac{N}{2(N-1)^2(N-3)},&N\text{ even},\\[6pt]
\dfrac{N+1}{2N^2(N-2)},&N\text{ odd}.
\end{cases}
\qquad\text{(22)}
\] Both are at most \(4/N^2\). For example, the joint crossing
probability before subtracting \(p^2\) is \[
\frac{4a(a-1)b(b-1)}
{N(N-1)(N-2)(N-3)}.
\]

There are \(\binom N2\) edges, \(3\binom N3\) unordered adjacent edge
pairs, and \(3\binom N4\) unordered disjoint edge pairs. Using absolute
values for the covariance terms, \begin{equation*}
\begin{aligned}
\operatorname{Var}Z
&\le L^2\left(
\frac14\binom N2
+\frac2N\,3\binom N3
+\frac8{N^2}\,3\binom N4
\right)\\
&\le \frac{17}{8}L^2N^2
<4L^2N^2.
\end{aligned}
\qquad\text{(23)}
\end{equation*}

Choose a basis \(u_1,\ldots,u_d\) of \(X\) with \[
\|u_j\|=1
\] and coordinate functionals \(\varphi_j\) of dual norm one. Such a
basis exists: maximize the absolute determinant of \(d\) vectors in the
unit ball. Replacing a basis vector by any other unit-ball vector shows
that every corresponding coordinate has absolute value at most one.
Also, \[
\|x\|\le\sum_{j=1}^d|\varphi_j(x)|.
\qquad\text{(24)}
\]

Set \[
T=\sum_e(I_e-p)h(e).
\] Applying (23) to the scalar weights \(\varphi_j(h(e))\), which have
absolute value at most \(L\), gives \begin{equation*}
\begin{aligned}
\mathbb E\|T\|
&\le \sum_{j=1}^d\mathbb E|\varphi_j(T)|\\
&\le \sum_{j=1}^d
\sqrt{\mathbb E[\varphi_j(T)^2]}\\
&\le 2LdN.
\end{aligned}
\end{equation*} Some partition therefore satisfies (20). \(\square\)

\paragraph{Completing the proof of Theorem
1}\label{completing-the-proof-of-theorem-1}

For \(N=3\), the Hamilton cycle consists of all three edges, and the
expression on the left of (2) is zero. Assume henceforth that \(N\ge4\).

Choose \(A,B\) as in Lemma 7, cyclically order \[
A=\{a_1,\ldots,a_a\},
\] and let \[
T=h(E(A,B)),\qquad S=h(K_N).
\]

If \(N\) is even, \(|B|=a\). If \(N\) is odd, adjoin a dummy symbol
\(\ast\) to \(B\), obtaining a set \(B^\ast\) of size \(a\).

Construct an auxiliary \(K_{a,a}\) between the cyclic gaps
\(g_i=(a_i,a_{i+1})\) and \(B^\ast\). For a real vertex \(v\in B\), set
\[
\lambda(g_i v)=h(a_i v)+h(a_{i+1}v).
\qquad\text{(25)}
\] In the odd case, set \[
\lambda(g_i\ast)=h(a_i a_{i+1}).
\qquad\text{(26)}
\] All auxiliary labels have norm at most \(2L\).

Every auxiliary perfect matching produces a Hamilton cycle: insert a
real vertex into its assigned gap, and in the odd case use the direct
edge across the dummy-assigned gap.

Let \(\mu\) be the mean auxiliary perfect-matching label sum. Then \[
\mu=
\begin{cases}
\dfrac2a T,&N\text{ even},\\[7pt]
\dfrac2a T+\dfrac1aY,&N\text{ odd},
\end{cases}
\qquad\text{(27)}
\] where, in the odd case, \[
Y=\sum_{i=1}^a h(a_i a_{i+1}).
\]

If \(N\) is even, then \[
\frac{2p}{a}=\frac2{N-1},
\] so \[
\mu-\frac2{N-1}S=\frac2a(T-pS).
\qquad\text{(28)}
\] If \(N\) is odd, then \(2p/a=2/N\), and consequently \[
\mu-\frac2{N-1}S
=
\frac2a(T-pS)
+\frac Ya-\frac{S}{\binom N2}.
\qquad\text{(29)}
\] Both \(Y/a\) and \(S/\binom N2\) have norm at most \(L\). Using Lemma
7 and \(a\ge N/2\), we obtain in either case \[
\left\|\mu-\frac2{N-1}S\right\|
\le 8Ld+2L.
\qquad\text{(30)}
\]

By Proposition 6, an auxiliary perfect matching has label sum within \[
6LdB_d
\] of \(\mu\). Its corresponding Hamilton cycle therefore satisfies (2).
This completes the proof of Theorem 1. \(\square\)

\begin{center}\rule{0.5\linewidth}{0.5pt}\end{center}

\subsubsection{7. Translation to colour
balance}\label{translation-to-colour-balance}

Give every edge of colour \(j\) the label \[
h(e)=e_j\in(\mathbb R^k,\|\cdot\|_1).
\] Thus every edge label has norm one, and \(h(F)\) is the colour-count
vector of \(F\).

For a globally colour-balanced \(K_{n,n}\), \[
\frac2n h(K_{n,n})=\frac{2n}{k}\mathbf1.
\] For a globally colour-balanced \(K_N\), \[
\frac2{N-1}h(K_N)=\frac Nk\mathbf1.
\]

If a subgraph \(F\) has \(kq\) edges and colour counts
\(m_1,\ldots,m_k\), its minimum recolouring number is \[
\rho(F)
=\sum_{j=1}^k(m_j-q)_+
=\frac12\sum_{j=1}^k|m_j-q|.
\qquad\text{(31)}
\] Indeed, every surplus edge must be recoloured, and recolouring
surplus edges into deficient colours attains this number.

Taking \(d=k\), \(L=1\), and halving the bounds in Theorem 1 proves (3)
and (4).

\subsubsection{Conclusion and scope}\label{conclusion-and-scope}

The proof establishes, uniformly in the host order, \[
\boxed{f_h(H)=O(k\log(k+1))}
\] in the bounded, zero-sum vector-label formulation, as well as the
explicit recolouring bounds above. It also establishes the corresponding
\(O(k\log(k+1))\) bipartite perfect-matching estimate.

Thus it answers the stated request for an improvement over \(O(k^2)\).
It does \textbf{not} establish an \(O(k)\) bound or determine the
optimal order. I have not independently checked whether a later revision
of the source already contains this improvement.

\appendix

\section{Provenance and review
status}\label{provenance-and-review-status}

\begin{warning}

\textbf{UNREFEREED MODEL OUTPUT.} This document was selected solely
because GPT-6 Astra labelled its own result
\texttt{would\_publish:\ true} and returned \texttt{proved} or
\texttt{disproved}. It has not passed the adversarial LLM referee used
for the earlier Sol campaign, has not been checked by a human
mathematician, and has not been checked for novelty. Treat every
mathematical and bibliographic claim as unverified.

\end{warning}

{\def\LTcaptype{none} % do not increment counter
\begin{longtable}[]{@{}
  >{\raggedright\arraybackslash}p{(\linewidth - 2\tabcolsep) * \real{0.5000}}
  >{\raggedright\arraybackslash}p{(\linewidth - 2\tabcolsep) * \real{0.5000}}@{}}
\toprule\noalign{}
\endhead
\bottomrule\noalign{}
\endlastfoot
Catalog id & \texttt{2604.09449\_\_03} \\
Catalog entry &
\href{https://graph-theory-ai.github.io/graph-conjectures/arxiv/2604.09449__03/}{Problem
6.4} \\
Source corpus & arxiv \\
Campaign leg & selected second attempts (\texttt{attacks\_retry}) \\
Source paper / entry & Colour-balanced subgraphs \\
Model verdict & \textbf{proved} (confidence: high) \\
Selection criterion & \texttt{would\_publish:\ true} and verdict
\texttt{proved} or \texttt{disproved} \\
Model & \texttt{gpt-6-astra}, reasoning effort \texttt{max},
\texttt{mode=pro}, flex service tier \\
Original artifact &
\texttt{attacks\_retry/2604.09449\_\_03/output.md} \\
Independent review & \textbf{None} \\
arXiv & \href{https://arxiv.org/abs/2604.09449}{arXiv:2604.09449} \\
Previous attempt & \texttt{attacks/2604.09449\_\_03}: gpt-5.6-sol
verdict \texttt{partial} \\
Model's caveats & Optimality and priority relative to later revisions of
the source paper are not established. \\
\end{longtable}
}

The generated Markdown and LaTeX sources for this PDF are stored under
\texttt{to\_review\_astra/src/2604.09449\_\_03/}.

\end{document}
